The UROP project aims to create a full stack,
web-based system that allows users to view the relationships between developers during the Code Review process.
Python is used to code the initial crawling stage of the project,
making use of the requests pip library
to make calls to the GitHub GraphQL API endpoint with an appropriate template string.
Python is also used in the migration stage of the project
to move the newly acquired pull request information from individual JSON files to an efficient relational database,
using the mysql-connector-python pip library to perform this.
The frontend handles the visualising of the data itself, and is primarily coded in TypeScript,
whilst the backend, coded in Java, makes queries to a local relational database to consummate the visualisation.
Axios comes in handy for the frontend stage of the project to make the call to the backend API endpoint,
and allows users to easily add parameters to a callback.
D3.js is an invaluable tool in deciphering the response Axios receives from that callback,
and subsequently produces the final visualisation the user sees.
Lastly,
Spring Boot was vital
in allowing easy development of the REST API building blocks that were essential to the development of the project.
This is particularly clear
when Spring Boot becomes
linked to the relational database
containing the pull request information since Spring Boot makes it incredibly simple
to make queries or to mutate the data inside the database tables,
if necessary.